\documentclass[UTF8,a4paper,11pt]{ctexart}
\usepackage[a4paper,left=25mm,right=25mm,top=25mm,bottom=25mm]{geometry}
\usepackage{amsmath,amssymb,bm}
\usepackage{booktabs}
\usepackage{graphicx}
\usepackage{xcolor}
\usepackage{tikz}
\usetikzlibrary{arrows.meta,positioning,calc,fit,shapes.geometric}
\usepackage[hidelinks]{hyperref}
\usepackage{enumitem}
\usepackage{ragged2e}
\usepackage{listings}
\usepackage[numbers,sort&compress]{natbib}

% 按中文学术论文习惯设置三级标题编号：一、/1.1/1.1.1。
\ctexset{
  section = {name={,、},number=\chinese{section},format=\zihao{-4}\bfseries,beforeskip=1.0em,afterskip=0.5em},
  subsection = {name={,},number=\thesection.\arabic{subsection},format=\zihao{5}\bfseries,beforeskip=0.8em,afterskip=0.35em},
  subsubsection = {name={,},number=\thesubsection.\arabic{subsubsection},format=\zihao{5}\bfseries,beforeskip=0.6em,afterskip=0.25em}
}
\setlength{\parindent}{2em}
\setlength{\parskip}{0.25em}
\setlength{\emergencystretch}{3em}
\hypersetup{pdftitle={面向 Android 应用测试的多源感知虚拟环境仿真框架设计与实现}}
\lstset{basicstyle=\ttfamily\small,breaklines=true,frame=single,columns=fullflexible}

\title{面向 Android 应用测试的多源感知虚拟环境仿真框架设计与实现}
\author{ZhangVirtualEnv 项目开发者}
\date{2026 年 8 月}

\begin{document}
\maketitle
\thispagestyle{plain}
\begin{center}
  {\zihao{-4}\bfseries 摘要}
\end{center}
\setlength{\parindent}{2em}
\setlength{\parskip}{0.25em}
\justifying
本项目面向 Android 应用开发、自动化测试和系统兼容性验证中的环境复现问题。实际开发过程中，单独修改位置坐标往往不能解释应用对卫星状态、NMEA、蜂窝小区、WiFi、BLE、SIM 和传感器数据的联合读取，测试结果因此难以复现，也难以判断问题究竟来自应用逻辑、Framework 接口还是厂商 ROM。针对这一问题，我在 ZhangVirtualEnv 中采用控制端、Backend Core、系统测试适配层和版本 Profile 分层的设计，将测试状态集中在统一的环境数据模型中，再根据 Android 版本和 ROM 差异把数据转换为 Framework 对象、Binder 返回值或传感器事件。本文以当前源码和 Android 15/Oplus 系统分析记录为基础，重点讨论路线位置的时间推进与插值、速度和方向计算、GNSS 卫星状态的几何生成、仰角与距离计算、C/N0 平滑模型，以及统一步态相位驱动的加速度、陀螺仪和计步器数据生成。同时，本文说明 CellInfo、WiFi、BLE 和 SIM 测试数据在系统服务边界的构造方式，并分析 SensorService Java 通道与 Native 事件分发通道之间的实际差异。现有实现已经形成单点和路线位置、GNSS 数据路径、CellInfo、WiFi、BLE、SIM 以及多级传感器测试后端；其中 Native 传感器通道是否在目标 ROM 上稳定送达，仍须通过独立检测器、Hook 日志和系统调用结果确认。已有验证表明，软件层测试链路能够提供连续、可观测的多源 API 数据，但它不等同于真实射频 GNSS、真实传感器硬件或物理无线环境。
\paragraph{关键词：} Android 环境测试；GNSS 模拟；系统服务适配；LSPosed；传感器仿真；位置服务；ROM 兼容性
\clearpage

\tableofcontents
\clearpage
\section{引言}
移动应用的行为越来越依赖设备所处的综合环境，而不仅是单个位置坐标。地图、运动、物联网和企业应用可能同时读取位置更新、卫星状态、NMEA、CellInfo、WiFi 扫描、BLE 广播、SIM 订阅信息以及加速度和计步事件。对应用开发者而言，真实外场测试成本高、环境难以复现，且 Android 不同版本和厂商 ROM 会改变类名、Binder 参数、服务拆分和事件分发路径。因此，需要一个能够保存测试数据、按时间回放并从系统接口边界提供可控返回值的测试框架。

仅修改\texttt{LocationManager} 或回调中的经纬度并不能形成一致测试环境。应用可以通过 \texttt{GnssStatus} 判断卫星是否可见，通过 NMEA 解析另一组坐标，通过 CellInfo/WiFi 估计网络位置，也可以通过传感器事件推断运动状态。若这些数据仍来自真实设备，应用内部的融合算法会得到相互矛盾的输入。本文将问题定义为：在不修改被测应用源码、也不向第三方应用进程安装测试适配的前提下，如何在 Android 系统服务与 Framework 边界提供可控、可回退、可观测的多源测试数据。

本文贡献如下：
\begin{enumerate}
  \item 总结 ZhangVirtualEnv 的分层架构和环境状态数据流，区分业务引擎、接口适配和版本 Profile 的职责。
  \item 从源码给出路线轨迹、GNSS 几何、信号强度和统一运动模型的数学描述。
  \item 结合 JADX 对 Android 15 \texttt{SensorService} 的结果，说明 Java Runtime Sensor 通道与 Native 分发通道的边界。
  \item 基于 VirEnvDetector 与既有记录，给出不夸大功能的验证结论和适用限制。
\end{enumerate}

本文将“模拟”限定为软件层测试数据生成与系统 API 适配，不将其表述为真实卫星射频、GNSS 芯片输入或物理网络环境重建。

\section{相关技术与问题分析}
\subsection{Android 感知数据通路}
Android 的位置接口由应用层 \texttt{LocationManager}、系统位置服务和 GNSS provider 共同完成；GNSS 状态、NMEA、导航消息和测量数据具有独立的注册与回调链~\citep{androidlocation}。蜂窝信息经 \texttt{TelephonyManager}、电话 Binder 服务和 \texttt{Phone} 对象层返回，CellIdentity 与 CellSignalStrength 共同构成可供上层解析的结构。WiFi 扫描结果由 WiFi 服务封装为 \texttt{ScanResult}，BLE 扫描则经蓝牙 Binder 与蓝牙栈回调 \texttt{ScanResult}。传感器事件通常经过 SensorService 的 native 事件队列，再进入应用的 \texttt{SensorEventListener}~\citep{androidsensor}。

\subsection{单一位置修改的不足}
设应用接收到的位置为 $L$，但同时观测到的卫星状态、无线环境和运动数据仍为真实值 $G,W,M$。应用的融合结果可表示为
\[
  \hat{P}=F(L,G,W,M).
\]
仅修改 $L$ 而保持 $G,W,M$ 不变，会使 $F$ 在不同数据源之间产生残差。例如，虚拟位置的 NMEA 坐标与真实卫星状态不一致，或虚拟速度与真实加速度、步频不一致。框架因此采用统一 Backend 状态，并让各适配器从同一测试 Profile 或时间轴读取数据。

\subsection{工程约束}
系统服务 Hook 具有高风险：错误的反射签名、错误的构造器字段或不兼容的 ROM 偏移可能影响 system\_server。项目采用候选类/方法匹配、版本 Profile、异常隔离和 fail-open 策略；适配失败时继续原始调用。scope.list 只包含必要系统组件、项目自身控制端和独立检测器，不把第三方应用作为实现依赖。

\section{系统总体架构设计}
我在实现过程中没有把各类 Hook 写成彼此独立的功能集合，而是先确定了一条统一的数据链路：控制端只负责提出测试场景，Backend Core 负责保存状态并推进时间，系统测试适配层负责把状态翻译成 Android 能够理解的接口数据。这样做的直接好处是，位置、卫星、无线环境和运动数据可以共享同一个时间基准，某一个 ROM 的类名或方法签名变化也不会反过来污染核心模型。系统结构如图~\ref{fig:architecture} 所示。

\begin{figure}[htbp]
\centering
\resizebox{0.88\linewidth}{!}{%
\begin{tikzpicture}[font=\small, node distance=10mm, box/.style={draw, rounded corners=2pt, align=center, minimum width=82mm, minimum height=11mm, fill=blue!7}, line/.style={-Latex, thick}]
  \node[box] (ui) {控制端 App\\测试参数、路线、Profile 与环境包};
  \node[box, below=of ui] (core) {Backend Core\\统一环境状态、时间轴、数据引擎与本地 API};
  \node[box, below=of core] (adapter) {系统测试适配层\\Location / GNSS / Cellular / WiFi / BLE / SIM / Sensor};
  \node[box, below=of adapter] (android) {Android Framework 与系统服务\\Binder、Java Framework、SensorService Native};
  \draw[line] (ui) -- node[right]{配置请求} (core);
  \draw[line] (core) -- node[right]{环境快照} (adapter);
  \draw[line] (adapter) -- node[right]{系统接口数据} (android);
\end{tikzpicture}}
\caption{系统总体架构与数据边界}
\label{fig:architecture}
\end{figure}

\subsection{状态与数据流}
配置请求沿“控制端 $\rightarrow$ 本地 API $\rightarrow$ Backend $\rightarrow$ 引擎”流动。适配器按周期或回调请求当前快照，生成具有当前时间戳的数据对象。位置是 GNSS 与部分环境生成器的共同参考，但 GNSS 引擎不持有第二份位置状态；路线引擎和单点引擎是位置的唯一所有者。环境采集器则把位置、CellInfo、WiFi、BLE、GNSS 和传感器事件按时间戳写入环境包，回放时由时间轴统一选择帧。

\subsection{版本 Profile}
Profile 按 SDK、设备指纹和 ROM 差异选择候选类名、方法名、参数数量和边界策略。Android 15 的蓝牙扫描入口与 Android 16 的 \texttt{ScanController.startScan} 存在差异，项目通过版本门控保持已验证版本行为不变。未匹配到目标类或方法时，适配器跳过该能力并放行系统原始路径，而不是以默认值覆盖未知 ROM。
\section{核心算法设计}
\subsection{虚拟轨迹模型}
我把位置生成分成单点和路线两种模式，但两者都只维护一份位置状态。单点模式保存经纬度、速度、方向和更新时间；路线模式保存折线点列、段索引和段内进度。对相邻点 $P_i=(\phi_i,\lambda_i)$ 与 $P_{i+1}$，先用 Haversine 近似计算段长
\[
 d_i=2R\arcsin\sqrt{\sin^2\frac{\Delta\phi}{2}+\cos\phi_i\cos\phi_{i+1}\sin^2\frac{\Delta\lambda}{2}}.
\]
在时间间隔 $\Delta t$ 内按 $v\Delta t$ 消耗路段长度；若未跨段，进度更新为 $q'=q+v\Delta t/d_i$，位置以线性插值给出
\[
 P(t)=(1-q)P_i+qP_{i+1}.
\]
路线引擎采用惰性推进：只有收到 \texttt{currentLocation()} 请求时才根据 monotonic 时钟推进，因此没有请求时不会由后台线程产生无观测数据。方向由相邻点的球面方位角计算，速度直接写入 \texttt{Location.speed}。循环路线可直接跳回起点，也可沿终点到起点的连线平滑返回。可选的小幅抖动只作用于输出坐标，不改变轨迹游标。

\subsection{GNSS 卫星状态模型}
GNSS 引擎以位置快照和时间戳为输入，对固定卫星目录逐颗计算状态：
\[
 S_k=\{\mathrm{SVID},\mathrm{constellation},\mathbf{r}_k, e_k, a_k, \rho_k, C/N_{0,k}, u_k\}.
\]
目录包含 GPS、BeiDou、Galileo、GLONASS 和 QZSS 测试卫星及主测试频段。轨道部分使用连续时间驱动的圆轨道 Kepler 近似。对模型参数 $\theta_k$，相位为
\[
 \alpha_k(t)=\alpha_{0,k}+2\pi (t\bmod T_k)/T_k,
\]
再按倾角和升交点将轨道平面坐标旋转到 ECEF。该算法提供 API 数据路径的连续几何，不包含星历、伪距、导航电文或射频波形。

\subsection{LLA、ECEF 与 ENU}
观测点的 WGS-84 LLA~\citep{wgs84} 为 $(\varphi,\lambda,h)$。令长半轴为 $a$、第一偏心率平方为 $e^2$，曲率半径
\[
 N=\frac{a}{\sqrt{1-e^2\sin^2\varphi}}.
\]
则
\[
 \mathbf{r}_{ecef}=\begin{bmatrix}(N+h)\cos\varphi\cos\lambda\\(N+h)\cos\varphi\sin\lambda\\(N(1-e^2)+h)\sin\varphi\end{bmatrix}.
\]
卫星与观测点的 ECEF 差向量经过局部 ENU 旋转得到 $(E,N,U)$。距离、仰角和方位角分别为
\[
 \rho=\sqrt{E^2+N^2+U^2},\quad e=\operatorname{atan2}(U,\sqrt{E^2+N^2}),
\]
\[
 a=\left(\operatorname{atan2}(E,N)\frac{180}{\pi}+360\right)\bmod 360.
\]
源码将仰角不低于 $5^\circ$ 的卫星视为可见，并要求仰角不低于 $15^\circ$ 且 $C/N_0\geq25$ dB-Hz 才标记为 UsedInFix。不可见行保留在诊断 JSON 中，但不伪造为公开 \texttt{GnssStatus} 行。

\subsection{C/N0 信号模型}
信号模型采用基准项、仰角增益、距离损耗和慢变化项：
\[
 C/N_0=\operatorname{clip}\left(C_0+G(e)-L(\rho)+n(t,k),18,52\right).
\]
其中 $G(e)$ 在低仰角为负、在高仰角趋于正增益，距离损耗为
\[
 L(\rho)=3\log_2\left(\max\left(\rho/r_k,0.5\right)\right),
\]
慢变化项由时间和卫星相位的正弦函数生成，而不是每次刷新独立抽样。这样可以保持 SVID、数量和几何连续，适用于解析器和状态机回归。

\subsection{统一运动与传感器模型}
\texttt{VirtualMotionEngine} 以单一 monotonic 时钟推进全局步态相位：
\[
 p_{t+\Delta t}=p_t+\frac{f\Delta t}{60}\pmod 1,
\]
其中 $f$ 为 steps/min。相位跨越 1 时，步数增加 1、距离增加一个步长，并产生一个待消费的 STEP\_DETECTOR 事件。四段步态波形模拟着地冲击、保持、摆动和恢复；加速度、重力、线性加速度、陀螺仪和计步器都由同一状态生成。手机姿态用低频 pitch/roll 漂移叠加步频摆臂，保证不同传感器之间具有时间相关性。Native C 实现复用了相同的相位、步数和姿态思想，并在传感器事件分发前重写数据。

\begin{figure}[htbp]
\centering
\resizebox{0.88\linewidth}{!}{%
\begin{tikzpicture}[font=\small, node distance=12mm and 15mm, box/.style={draw, rounded corners=2pt, align=center, minimum width=36mm, minimum height=11mm, fill=orange!10}, line/.style={-Latex, thick}]
\node[box] (pos) {统一 Location 快照\\路线/单点引擎};
\node[box, right=of pos] (orbit) {卫星目录与时间\\圆轨道近似};
\node[box, below=of pos] (coord) {LLA $\rightarrow$ ECEF\\ECEF $\rightarrow$ ENU};
\node[box, right=of coord] (signal) {仰角、距离\\可见性、C/N0};
\node[box, below=of coord] (status) {GnssStatus / NMEA\\诊断 JSON};
\draw[line] (pos) -- (coord);
\draw[line] (orbit) -- (signal);
\draw[line] (coord) -- (signal);
\draw[line] (signal) -- (status);
\end{tikzpicture}}
\caption{GNSS 状态生成过程}
\label{fig:gnss}
\end{figure}
\section{Android 实现方法}
\subsection{Location 与 GNSS 适配}
Location 适配器在系统位置服务边界接收注册、上报和监听器相关调用，并从 Backend 取得当前测试位置。GNSS 适配器覆盖状态和 NMEA 注册及注销路径；启用测试位置时，不注册真实回调，而是周期投递由同一位置快照生成的 \texttt{GnssStatus} 和 NMEA。导航消息与原始测量属于更低层数据面，当前实现对其采用阻断或预留策略，关闭测试时恢复原始路径。\texttt{GnssLocationProvider.onReportSvStatus(GnssStatus)} 与 \texttt{onReportNmea(long,String)} 同时用于观察真实系统数据和替换测试数据，所有异常均回退到原始调用。

\subsection{Cellular 与 SIM}
CellInfo 测试数据按 LTE、NR、GSM、WCDMA 和 CDMA 组织。构造器使用 Profile 提供的 MCC/MNC、TAC/LAC、CI/NCI、PCI、ARFCN 和信号强度，并对字段位宽和不可用哨兵进行归一化。Android 15 的 MCC/MNC 字符串字段与旧 int getter 存在差异，读取器优先使用字符串并安全解析。CDMA 位置字段需要按定点量化方式写入 CellIdentity，故其是特定 ROM 的兼容构造路径，而不是通用蜂窝定位模型。

SIM 适配位于电话 Binder 服务端和 Phone 对象层，覆盖订阅、运营商、设备标识、号码、网络类型、SIM 状态和信号强度等返回值。实现对无后缀、ForSubscriber、WithFeature 等方法名建立候选集合；目标不存在或返回值构造失败时调用原始方法。该设计强调接口抽象和兼容性测试，不表示对真实基带或 SIM 卡物理状态的重建。

\subsection{WiFi 与 Bluetooth}
WiFi 服务由 ServiceManager 动态发现，以应对 WiFi APEX 中类不在 boot classloader 的情况。\texttt{getScanResults} 返回由 SSID、BSSID、RSSI、频率和 capabilities 构造的 \texttt{ScanResult} 列表；\texttt{getConnectionInfo} 返回对应连接信息。对 Oplus 扩展字段如 information elements 做非空初始化，避免系统消费者对空数组解引用。BLE 测试数据按地址去重，使用 \texttt{ScanResult(BluetoothDevice, ScanRecord, rssi, timestamp)} 构造广播结果，原始广播字节优先，否则按名称生成最小广播记录。蓝牙栈 Hook 根据 Android 15/16 的扫描入口差异选择候选方法，并在无法解析回调时放行真实扫描。

\subsection{Sensor Native 适配}
JADX 结果显示目标 ROM 的传感器 Java 外壳为 \texttt{com.android.server.sensors.SensorService}，持有 native 指针 \texttt{mPtr}，通过 \texttt{SensorManagerInternal.LocalService} 暴露 Runtime Sensor 的创建、删除和事件发送接口。其 \texttt{sendRuntimeSensorEventNative(long,int,int,long,float[])} 只对服务注册的 runtime handle 有效；对真实传感器 handle 调用成功并不等价于事件已进入普通应用连接。项目已有记录据此将 Java 通道作为探测/回退路径，并把全局研究重点放在 \texttt{libsensorservice.so} 的分发汇聚点。

Native 方案在编译期按目标 ROM 的函数特征校验入口，在 \texttt{SensorService::sendEventsToAllClients} 处读取服务事件缓冲，重写事件值并尾跳原函数；\texttt{SensorEventQueue::write} 和 \texttt{BitTube::sendObjects} 是保留的兜底入口。该路径需要严格的 ARM64 指令、PAC、内存保护和容量边界处理，安装失败时不修改原代码。论文将其作为系统级测试适配的实现路径；其是否对未配置作用域的测试进程稳定送达，仍应以 VirEnvDetector 的事件计数和系统日志进行运行时确认。

\begin{figure}[htbp]
\centering
\resizebox{0.88\linewidth}{!}{%
\begin{tikzpicture}[font=\small, node distance=11mm and 17mm, box/.style={draw, rounded corners=2pt, align=center, minimum width=39mm, minimum height=11mm, fill=purple!8}, line/.style={-Latex, thick}]
\node[box] (cfg) {Profile 与本地 API\\测试参数};
\node[box, below=of cfg] (backend) {Backend\\Environment State};
\node[box, below left=12mm and 2mm of backend] (java) {Java 适配器\\Binder / Framework};
\node[box, below right=12mm and 2mm of backend] (native) {Native 适配器\\Sensor 分发边界};
\node[box] (app) at ($(java.south)!0.5!(native.south)+(0,-18mm)$) {测试目标通过正常 Android API\\接收测试数据};
\draw[line] (cfg) -- (backend);
\draw[line] (backend) -- (java);
\draw[line] (backend) -- (native);
\draw[line] (java.south) -- (app.north west);
\draw[line] (native.south) -- (app.north east);
\end{tikzpicture}}
\caption{测试数据的系统适配流程}
\label{fig:dataflow}
\end{figure}
\section{实验与结果}
\subsection{实验方法}
本文不重新执行真机测试，而复用项目已有的源码、逆向记录和 VirEnvDetector~\citep{virenvdetector} 设计作为实验材料。验证原则是独立检测器优先，并结合模块 Hook 状态、logcat、系统调用结果和检测器报告；截图不作为主要证据。检测器从普通应用视角注册位置、GNSS Status/NMEA、传感器、BLE、WiFi 和电话读取，并通过本地 API 获取期望配置，形成 PASS、FAIL、SYNCING 或 NOT\_ENABLED 判定。结果解释区分“接口回调到达”“数据字段一致”和“Native 全局通道实际送达”三个层级。

\subsection{源码与静态证据}
源码确认以下能力已形成明确实现链路：单点位置由 \texttt{SinglePointLocationEngine} 管理，路线由惰性时间推进的 \texttt{RouteEngine} 管理；GNSS 由 \texttt{GNSSSimulationEngine}、\texttt{OrbitCalculator}、\texttt{CoordinateTransform} 和 \texttt{SignalModel} 组合生成；CellInfo、WiFi、BLE、SIM 由对应工厂和系统服务适配器构造。运动数据由 \texttt{VirtualMotionEngine} 统一生成，传感器注入层仅负责投递~\citep{zhangvirtualenv}.

\subsection{已有运行时结论}
项目记录显示，目标 Android 15/Oplus 设备上的应用兼容传感器链路能够使 VirEnvDetector 观察到步数递增，以及由统一运动模型产生的加速度、重力和磁场测试值；注册真实传感器的路径在模拟激活时被屏蔽，避免测试值与真实事件叠加。GNSS 记录显示，虚拟状态包含卫星数量、星座、SVID、角度、CN0、星历和 UsedInFix 字段，且 NMEA 与虚拟位置使用同一 Backend 位置源。

对系统级 SensorService Java 通道，已有记录明确指出 Oplus 15 的 Runtime Sensor 事件调用可以返回成功并更新缓存，但对真实传感器 handle 未进入普通应用连接；因此不能将该返回值当作全局送达证据。Native 汇聚点方案提供了更接近真实分发路径的实现，但必须用 detector 的 \texttt{events received} 增长、\texttt{dumpsys sensorservice}、Hook 计数和 logcat 共同确认。该证据边界使论文不把设计预留写成已经完成的全局能力。

\subsection{兼容性结论}
兼容矩阵将 Android 15 多项能力标记为 VERIFIED，将 Android 16 多项类和方法签名标记为 STATIC\_VERIFIED，并将 WiFi 动态类、前台服务规则及运行时 SIM 链路保留为 REQUIRES\_DEVICE。Android 17 及以上为 NOT\_ADAPTED。由此，当前结论是：在已验证的软件测试环境中表现稳定；跨版本可通过 Profile 维持候选匹配和回退，但不能由静态签名推导真机运行成功。

\section{讨论与限制}
\subsection{软件层边界}
框架生成的是 Android API 可观察的数据，不是 GNSS 射频信号。GNSS 轨道采用圆轨道近似，未模拟真实星历、伪距、载波相位、多路径、接收机时钟误差、导航电文和射频前端；因此不能替代硬件在环或射频仿真器。类似地，WiFi、BLE 和 CellInfo 是结构化测试对象，不会改变物理网络、无线电环境或基带状态。

\subsection{厂商 ROM 差异}
系统服务可能移动到 APEX、厂商 JAR 或独立 APK，R8 可能重命名内部方法，Binder 参数也可能随版本变化。Native Hook 还依赖函数布局、指令特征、PAC 和可执行内存权限。项目通过 Profile、多候选匹配、二进制特征检查和 fail-open 降低风险，但这不能保证未知 ROM 兼容。任何涉及 system\_server 的新适配都应先静态分析，再在独立检测器上做最小运行验证。

\subsection{数据一致性与可观测性}
环境多源一致性依赖统一时间轴。若适配器缓存滞后、监听器生命周期处理错误或某一数据面继续返回真实值，融合应用仍可能观察到残差。因此测试报告应记录配置版本、快照时间、Hook 状态、回调时间、原始/虚拟摘要和系统日志，而不能只记录 UI 文本。对于事件驱动计步器，静止状态没有真实事件，必须区分“没有事件”与“值为零”，并在需要时使用主动事件机制。

\subsection{后续方向}
后续工作包括：为 Android 14、16 及后续版本补充可复现实机证据；建立不依赖具体 ROM 偏移的 Sensor HAL 或系统扩展适配；将路线速度与统一运动状态双向约束；加入更严格的时间戳、单调性和跨源一致性指标；并把环境包回放与 CI 测试用例结合。上述方向均属于测试基础设施增强，不改变项目不针对第三方应用进程开发的架构约束。

\section{总结}
本文从问题背景、系统架构、数学模型、Android 实现和实验边界五个方面分析了 ZhangVirtualEnv。框架以 Backend Core 维护统一环境状态，以系统服务和 Framework 边界的适配器将状态转换为 Android API 数据，并通过 Profile 隔离 Android 版本与厂商 ROM 差异。路线引擎实现按时间推进的单点和折线轨迹，GNSS 引擎实现固定卫星目录、连续轨道近似、WGS-84 坐标变换、ENU 几何和连续 C/N0；运动引擎以统一步态相位生成步数、加速度、重力和陀螺仪数据；CellInfo、WiFi、BLE 和 SIM 适配器则通过版本兼容构造器和 Binder 服务边界提供结构化测试输入。

源码和逆向证据同时表明，数据路径的“可调用”并不等于“已送达”。尤其是 Android 15 Oplus SensorService 的 Runtime Sensor 通道存在真实传感器句柄边界，Native 分发汇聚点虽是合理的系统级研究方向，也必须由 VirEnvDetector、Hook 计数、logcat 与系统调用共同证明。因此，本文最终将 ZhangVirtualEnv 定位为可重复的 Android 系统环境参数测试与兼容性验证框架，而非真实物理环境或射频信号仿真器。


\clearpage
\bibliographystyle{plainnat}
\bibliography{references}
\end{document}
